\section{Kết luận}

Đề tài đã xây dựng được một nguyên mẫu hệ thống IoT giám sát chất lượng không khí trong
phòng học, bao gồm các thành phần chính: nút ESP32 thu thập dữ liệu cảm biến, firmware
RTOS xử lý cục bộ, kênh truyền thông MQTT, backend Node.js/Express, cơ sở dữ liệu MySQL
và dashboard web phục vụ giám sát, phân tích và điều khiển. Kết quả này cho thấy kiến
trúc tổng thể của hệ thống đã được hình thành đầy đủ ở mức nguyên mẫu và các tầng có thể
phối hợp với nhau theo một contract dữ liệu thống nhất.

Ở tầng thiết bị, hệ thống đã thu thập được các nhóm đại lượng môi trường quan trọng gồm
CO\textsubscript{2}, bụi mịn, nhiệt độ, độ ẩm và VOC. Dữ liệu được chuẩn hóa về cùng mô
hình, đánh giá bằng cơ chế ngưỡng có hysteresis và ánh xạ sang trạng thái cảnh báo cũng
như trạng thái thiết bị xử lý cục bộ. Nhờ đó, thiết bị không chỉ thực hiện chức năng đo
lường mà còn có khả năng phản ứng tại chỗ khi điều kiện môi trường vượt khỏi vùng vận
hành mong muốn.

Ở tầng hệ thống, firmware RTOS đã được tổ chức thành các tác vụ tương đối độc lập cho
đọc cảm biến, xử lý IAQ, điều khiển, cảnh báo và truyền thông MQTT. Tầng backend tiếp
nhận telemetry, lưu lịch sử, cung cấp API và hỗ trợ luồng điều khiển ngược từ dashboard
đến ESP32. Trên cơ sở đó, dashboard web đã cung cấp được ba nhóm chức năng cốt lõi:
quan sát realtime, xem lại dữ liệu lịch sử và thực hiện điều khiển hoặc cấu hình từ xa.

Nhìn chung, đề tài đã chứng minh được tính khả thi của pipeline
\texttt{Sensor -> ESP32/RTOS -> MQTT -> Backend -> Database -> Dashboard}. Đây là kết quả
quan trọng đối với phạm vi đồ án, vì hệ thống không dừng ở mức đọc cảm biến đơn lẻ mà đã
đạt tới một mô hình IoT hoàn chỉnh ở mức nguyên mẫu, có dữ liệu thống nhất, có lưu trữ,
có giao diện và có khả năng điều khiển hai chiều.

\section{Giới hạn hiện tại}

Mặc dù đã đạt được mục tiêu ở mức nguyên mẫu, hệ thống hiện tại vẫn còn một số giới hạn.
Trước hết, các đầu ra điều khiển mới dừng ở mức mô phỏng bằng LED, chưa thay thế bằng
thiết bị chấp hành công suất thực như quạt thông gió, máy lọc không khí hoặc relay điều
khiển tải. Vì vậy, các vấn đề về an toàn điện, cách ly công suất và vận hành thiết bị
thực chưa nằm trong phạm vi hoàn thiện của phiên bản hiện tại.

Thứ hai, chỉ số VOC đang được sử dụng dưới dạng raw/index tương đối, chưa được hiệu
chuẩn thành đại lượng định lượng có ý nghĩa tham chiếu rộng hơn. Tương tự, các ngưỡng
điều khiển hiện nay phù hợp cho mục tiêu kiểm thử và minh họa cơ chế vận hành, nhưng chưa
thể xem là bộ ngưỡng cuối cùng cho triển khai thực địa nếu chưa có quá trình đo đạc, đối
chiếu và hiệu chuẩn trong môi trường sử dụng thực tế.

Thứ ba, hệ thống mới được kiểm chứng tốt ở mức nguyên mẫu kỹ thuật và các bài kiểm thử
thành phần. Dù backend đã có test cho MQTT và API, hệ thống vẫn cần được đánh giá
end-to-end trong điều kiện vận hành liên tục với ESP32 thực, broker thực, cơ sở dữ liệu
và dashboard hoạt động đồng thời trong thời gian dài. Đây là bước cần thiết để đánh giá
độ ổn định, độ trễ, khả năng phục hồi kết nối và tỷ lệ mẫu đo hợp lệ.

Cuối cùng, các yêu cầu về bảo mật và quản trị vận hành hiện vẫn ở mức cơ bản. MQTT và API
chủ yếu phục vụ môi trường lab hoặc demo; các cơ chế như xác thực người dùng, phân quyền,
mã hóa kết nối và audit log vẫn cần được bổ sung nếu hệ thống được mở rộng sang môi
trường triển khai thực tế.

\section{Hướng phát triển}

Hướng phát triển đầu tiên là hoàn thiện tầng phần cứng và thiết bị chấp hành. Các LED mô
phỏng có thể được thay bằng relay, MOSFET hoặc các module điều khiển phù hợp với quạt
thông gió, thiết bị lọc bụi, thiết bị lọc than, điều hòa hoặc máy tạo ẩm. Khi đó, hệ
thống sẽ tiến gần hơn tới một mô hình điều khiển IAQ có khả năng ứng dụng thực tế.

Hướng thứ hai là hiệu chuẩn cảm biến và tinh chỉnh bộ ngưỡng vận hành. Việc đối chiếu
dữ liệu với thiết bị tham chiếu, thu mẫu trong nhiều điều kiện lớp học khác nhau và đánh
giá lại ngưỡng CO\textsubscript{2}, PM2.5, VOC, nhiệt độ và độ ẩm sẽ giúp hệ thống có cơ
sở vận hành chặt chẽ hơn. Đối với VOC, đây cũng là bước cần thiết để chuyển từ raw/index
sang một chỉ số có ý nghĩa ứng dụng cao hơn.

Hướng thứ ba là tăng độ tin cậy và khả năng vận hành dài hạn của firmware RTOS. Các cải
tiến có thể bao gồm watchdog, giám sát tài nguyên của từng task, cơ chế lưu đệm dữ liệu
khi mất kết nối MQTT và đồng bộ bù khi mạng phục hồi. Những bổ sung này sẽ giúp hệ thống
ổn định hơn khi chuyển từ nguyên mẫu sang vận hành liên tục.

Hướng thứ tư là mở rộng backend và dashboard theo nhu cầu sử dụng thực tế. Các khả năng
có thể được bổ sung gồm cập nhật realtime bằng WebSocket, thống kê dài hạn, cảnh báo qua
email hoặc ứng dụng nhắn tin, và quản lý nhiều thiết bị theo \texttt{device\_id}. Khi số
lượng thiết bị tăng lên, đây sẽ là điều kiện cần để hệ thống chuyển từ mô hình demo đơn
thiết bị sang mô hình giám sát phân tán.

Cuối cùng, một hướng phát triển quan trọng là hoàn thiện lớp bảo mật và vận hành. Việc
bổ sung xác thực cho API, bảo vệ broker MQTT, quản lý cấu hình thiết bị và hỗ trợ cập
nhật firmware sẽ giúp hệ thống phù hợp hơn với yêu cầu triển khai ngoài môi trường phòng
thí nghiệm.
